\chapter{Conclusion and Future Work}
\label{ch:conclusion}

%% ── 8.1 ─────────────────────────────────────────────────────────────────────
\section{Summary}
\label{sec:summary}

% Restate: problem → approach → results

% "This thesis presented a federated TinyML framework for environment-driven
%  path loss regression on LoRaWAN edge devices. Building on a prior eight-month
%  indoor measurement campaign that established the relationship between
%  environmental factors and LoRaWAN path loss (centralized XGBoost R² ≈ 0.93),
%  this work addressed three key limitations of the centralized approach:
%  excessive bandwidth consumption, lack of data privacy, and inability to
%  adapt post-deployment.
%
%  The proposed system deploys a compact neural network (57 parameters, 2.7 KB)
%  on Arduino MKR WAN 1310 nodes. Each device performs local inference and
%  training, transmitting only quantized weight updates (52 bytes) instead of
%  raw sensor data (18 bytes × 1440/day). Server-side FedAvg aggregation
%  combines updates from multiple devices into an improved global model.
%
%  Evaluation on 1,715,869 real measurements, partitioned by physical device
%  location to simulate realistic non-IID conditions, demonstrates that the
%  federated approach achieves R² = [XX], compared to centralized R² = [XX]—
%  a reduction of only [X]%—while reducing daily uplink bandwidth by
%  approximately 11×. Hardware deployment confirms the pipeline fits within
%  the MKR WAN 1310's 32 KB SRAM and 256 KB Flash."

%% ── 8.2 ─────────────────────────────────────────────────────────────────────
\section{Key Findings}
\label{sec:key-findings}

\begin{enumerate}
  \item \textbf{Federated learning achieves comparable accuracy.}
    $R^2 = $ [XX] vs.\ centralized $R^2 = $ [XX], a difference of only [X]\%.

  \item \textbf{Communication reduced by $\approx$11$\times$.}
    From $\approx$\SI{25.9}{\kilo\byte}/day/node to $\approx$\SI{2.4}{\kilo\byte}/day/node.

  \item \textbf{Fits on real constrained hardware.}
    [XX]\% Flash, [XX]\% SRAM utilization on the MKR~WAN~1310.

  \item \textbf{Single-message FL updates.}
    Unlike Torres Sanchez et al.'s 1.39\,KB model requiring 7+ messages, our
    52-byte update fits in a single \gls{lorawan} uplink at any spreading factor.

  \item \textbf{Non-\gls{iid} impact is [manageable/significant].}
    Device-based data partitioning [did/did not] significantly impair convergence
    compared to \gls{iid} splits.
\end{enumerate}

%% ── 8.3 ─────────────────────────────────────────────────────────────────────
\section{Future Work}
\label{sec:future-work}

\begin{enumerate}
  \item \textbf{Live federated deployment.}
    Deploy all six nodes for a multi-week live campaign where devices perform real
    federated rounds through actual LoRaWAN communication, validating simulation
    results in the field.

  \item \textbf{True on-device backpropagation.}
    Replace the simplified weight proxy with actual gradient computation on the
    TFLite model when TFLite Micro adds training support (TinyTrain, TinyOL).

  \item \textbf{TTN downlink API integration.}
    Complete the server-side downlink scheduling to automatically push global
    model updates to devices via the TTN~v3 REST~API.

  \item \textbf{Architecture exploration.}
    Evaluate deeper networks (e.g., Dense~5$\to$16$\to$8$\to$1) and assess
    the accuracy--communication trade-off for larger models.

  \item \textbf{Energy profiling.}
    Measure the battery impact of local training cycles on the MKR~WAN~1310
    to estimate operational lifetime under federated workloads.

  \item \textbf{Cross-environment generalization.}
    Test whether a model trained in one building can transfer to a different
    building via federated fine-tuning.
\end{enumerate}
